\documentclass[11pt]{article}
\usepackage[a4paper,margin=1in]{geometry}
\usepackage{amsmath,amssymb}
\usepackage{booktabs}
\usepackage{graphicx}
\usepackage{siunitx}
\usepackage[hidelinks]{hyperref}
\usepackage{xcolor}
\usepackage[section]{placeins}
\graphicspath{{../figures/}}

\title{\bfseries Multi-Curve Calibration of a Two-Factor Lifted-CIR\\
Hull--White--Heston Model to Caplet Volatility Surfaces}
\author{Roland Grinis\thanks{\href{mailto:roland.grinis@grinisrit.com}{\texttt{roland.grinis@grinisrit.com}}}
\and Andrey Popov}
\date{\today}

\newcommand{\dW}{\,\mathrm{d}W}
\newcommand{\dt}{\,\mathrm{d}t}
\newcommand{\bp}{\,\text{bp}}
\newcommand{\doi}[1]{\href{https://doi.org/#1}{doi:\nolinkurl{#1}}}

\begin{document}
\maketitle

\begin{abstract}
We calibrate a multi-curve, two-factor \emph{lifted}-CIR stochastic-volatility short-rate model
(a Hull--White--Heston, HWH, construction) jointly to the entire EUR caplet normal-volatility
surface. The state is driven by a shared level factor with stochastic variance and a Hull--White
basis factor separating the projection curve from the OIS discount curve. Caplets are priced by
GPU Monte Carlo on the instrument each market actually trades---a forward-looking Euribor term
rate in EUR, a compounded in-arrears overnight rate in USD and GBP---and inverted to Bachelier
(normal) volatilities on the matching expiry convention. The CIR variance factors are sampled with the Andersen (2008) quadratic-exponential
(QE) scheme, and the 27 parameters are estimated with a gradient-free Adam--SPSA optimiser under a
vega-weighted objective. Against 340 caplets spanning maturities \(0.5\)--\(10\)Y and strikes
\(-1.5\%\) to \(+10\%\), the model attains an overall normal-volatility RMSE of \(5.9\bp\), an
at-the-money (ATM) mean absolute error of \(2.0\bp\), and a \(2.8\bp\) RMSE over the liquid
(\(|K-F|<2\%\)) core. Two corrections matter as much as the model: pricing the Euribor caplet as the
forward-looking contract it is, rather than as an in-arrears average, cuts the EUR error from
\(8.6\) to \(6.3\bp\), and a deterministic shift that makes the model reprice the initial curves
cuts the liquid-core error from \(3.7\) to \(2.8\bp\). Simulation is on a \emph{daily} grid (\(dt=1/252\)), on which every accrual
window is a whole number of steps. A nested-model ladder shows that stochastic volatility is
responsible for the bulk of the fit: the best flat-smile Gaussian restriction attains only
\(\approx 45\bp\). The same model reproduces the USD (SOFR) and GBP (SONIA) in-arrears caplet
surfaces from \(6\)M maturity out to \(8.0\) and \(6.7\bp\), with ATM errors of \(2.2\) and \(4.2\bp\).
\end{abstract}

\medskip
\noindent\textbf{Keywords:} multi-curve interest rates; stochastic volatility; lifted Heston;
caplet/cap calibration; normal (Bachelier) implied volatility; Monte Carlo; SPSA.

\smallskip
\noindent\textbf{JEL classification:} C63; G12; G13.

\section{Motivation}
EUR caplet surfaces exhibit a pronounced, maturity-dependent smile: a
convex ``payer'' wing at high strikes and a steep short-maturity slope that flattens with maturity.
Any Gaussian short-rate model~\cite{vasicek1977,hullwhite1990,brigo2006,filipovic2009} (one- or multi-factor
Hull--White with a deterministic volatility term structure) produces a \emph{strike-flat}
normal-volatility smile and therefore cannot, by construction, reproduce this curvature. Introducing
an affine stochastic variance~\cite{cir1985,heston1993,duffie2000} for the driving
factor---as in the SABR~\cite{hagan2002} and stochastic-volatility literature---restores both the
curvature (through the vol-of-vol) and the skew (through a leverage
correlation), while keeping the model tractable for large-scale Monte Carlo calibration. The
purpose of this note is to document precisely \emph{what} process is fitted, \emph{how} the market
data are fed in, and \emph{how well} the joint calibration reproduces the full surface.

\paragraph{Contributions.} The construction combines six ingredients into one calibratable object.
(i) A genuinely \emph{multi-curve} state (OIS discount + Euribor projection, coupled through a
shared factor and a basis factor), within the multiple-yield-curve framework of Cuchiero, Fontana
and Gnoatto~\cite{fontana2016hjm,cfg2019}, so the smile is priced consistently across both curves.
(ii) A \emph{lifted} stochastic variance---a two-factor CIR mixture that is a low-order Markovian
lift, in the sense of Abi~Jaber and El~Euch~\cite{abijaber2019multi} and
Cuchiero and Teichmann~\cite{cuchiero2019}, of a rough/multi-scale volatility---built from affine (CIR) factors and
directly simulable. (iii) A \emph{large-scale joint calibration}: one \(27\)-parameter vector fit to
the entire caplet surface by GPU Monte Carlo and a gradient-free Adam--SPSA
optimiser~\cite{spall1998,kingma2015}. (iv) The Andersen QE scheme~\cite{andersen2008} for the
Feller-violated regime (\(2\kappa\theta\ll\varepsilon^2\)) in which the calibrated variance factors
live. (v) A vega-weighted objective that targets the liquid core. (vi) Caplets priced as the contract
each market actually trades---a forward-looking Euribor term rate in EUR, a compounded in-arrears
overnight rate in USD and GBP~\cite{lyashenko2019,fontana2023caplet}---with the model repricing the
initial curves exactly. We report the EUR fit and its robustness across USD and GBP.

\section{Related literature}
\label{sec:lit}
\paragraph{Stochastic volatility in term-structure models.} Gaussian short-rate
models~\cite{vasicek1977,hullwhite1990} give strike-flat normal smiles.
Collin-Dufresne and Goldstein~\cite{collindufresne2002} showed that cap and swaption volatility is
driven by factors that bonds do not span---the empirical case for a separate variance factor, which
our nested ladder (Section~\ref{sec:ladder}) confirms on the caplet surface. Trolle and
Schwartz~\cite{trolleschwartz2009} calibrate a multi-factor HJM model with unspanned stochastic
volatility to caps and swaptions, and quasi-Gaussian models with a CIR variance multiplier are the
practitioner standard~\cite{andersenpiterbarg2010}. Our model sits in
this family: a Hull--White level factor whose variance is a two-factor CIR mixture. The alternative
is to model the quoted forwards directly, with the LIBOR market model~\cite{bgm1997} and the SABR
smile~\cite{hagan2002}, which fit slice by slice; a short-rate model fits the whole surface with one
parameter vector and gives every discount factor at every date, which is what Monte Carlo pricing of
in-arrears payoffs needs.

\paragraph{Affine term-structure models.} Affine short-rate models go back to Duffie and
Kan~\cite{duffiekan1996}; Duffie, Filipovi\'{c} and Schachermayer~\cite{dfs2003} characterise affine
processes on the canonical state space \(\mathbb{R}_+^m\times\mathbb{R}^n\), Duffie, Pan and
Singleton~\cite{duffie2000} give the transform analysis that makes them tractable. Our state \((x,V_1,V_2)\) lives on
\(\mathbb{R}\times\mathbb{R}_+^2\) and is affine apart from one term: the leverage covariance
\(\rho\,\varepsilon\sqrt{V V_1}\) is not linear in the state unless \(w=1\).

\paragraph{Lifted volatility.} Volatility is rough~\cite{gatheral2018}, and its finite-dimensional
Markovian lifts~\cite{abijaber2019multi,abijaber2019lifting,cuchiero2019} replace the fractional
kernel by a weighted sum of mean-reverting factors; our two-factor CIR mixture is the lowest-order
such lift, applied to a rate rather than an equity variance.

\paragraph{Multi-curve models and backward-looking rates.} The separation of discounting from forward
estimation became standard after 2008~\cite{bianchetti2010,mercurio2010,henrard2026}. Models with a
stochastic basis include Mercurio and Xie~\cite{mercurioxie2012}. Cuchiero, Fontana and Gnoatto set out a general HJM framework for
multiple curves~\cite{fontana2016hjm} and classify its affine specifications~\cite{cfg2019}, and our
Hull--White basis factor is the simplest member of that class; Fontana, Gnoatto and
Szulda~\cite{fontana2021cbi} drive the spreads by self-exciting CBI processes, Fontana, Lanaro and
Murgoci~\cite{fontana2025geometry} characterise which multi-curve HJM models admit
finite-dimensional factor realisations, and Fontana, Pavarana and Runggaldier~\cite{fontana2023rollover}
give the Euribor--OIS basis an economic reading as roll-over risk. Since the
LIBOR transition, USD and GBP caps pay a compounded overnight rate that is known only at the end of
the accrual period, while EUR caps still reference the forward-looking Euribor. Lyashenko and
Mercurio~\cite{lyashenko2019} extended the market model to such in-arrears rates, and
Fontana~\cite{fontana2023caplet} prices them in general affine short-rate models---the setting closest to ours. Overnight rates jump at policy
dates: Fontana, Grbac, G\"umbel and Schmidt~\cite{fontana2020disc} and Fontana, Grbac and
Schmidt~\cite{fontana2024overnight} build term-structure models with such stochastic discontinuities,
and Andersen and Bang~\cite{andersenbang2024} apply jumps to SOFR caplets. The common point, which our instrument definition
(Section~\ref{sec:instrument}) follows, is that a backward-looking caplet keeps volatility through
its accrual period and a forward-looking one does not.

\section{The multi-curve model}

\subsection{Why two curves}
Since 2008 the single-curve identity---interbank term rates (Euribor) treated as risk-free and equal
to OIS---no longer holds: Euribor carries a credit/liquidity premium, while collateralised
derivatives are discounted at the collateral rate, the overnight rate (€STR). Pricing a EUR cap therefore needs
two curves: a \textbf{discount} curve (€STR/OIS) for the collateralised payoff, and a
\textbf{projection} curve (6M Euribor) for the floating rate it pays on. In the EUR data the two
differ by \(+9\) to \(+21\bp\) (the Euribor--OIS basis); collapsing them mis-prices every forward by
that amount---an order of magnitude above our few-bp fit target. This is the multi-curve framework
of Bianchetti~\cite{bianchetti2010}, Mercurio~\cite{mercurio2010}, Henrard~\cite{henrard2026}
and others~\cite{grbac2015}.
Affine specifications of that framework are classified by Cuchiero, Fontana and
Gnoatto~\cite{cfg2019}, and the state process below sits inside their affine multiple-yield-curve
class; our contribution is to pair that structure with a lifted-CIR stochastic variance and to
calibrate the result, so the \emph{smile}, not just the forward, is priced consistently across
both curves.

\subsection{State dynamics}
Let \(f_{\text{ois}}(t)\) and \(f_{\text{proj}}(t)\) denote the OIS (discount) and 6M-Euribor
(projection) instantaneous forward curves, both taken directly from the market. We drive them with
a \emph{shared} level factor \(x\)---the common, essentially risk-free, movement of all EUR
rates---and an additive Hull--White \emph{basis} factor \(k\) that lives only on the projection
curve:
\begin{align}
r_{\text{ois}}(t) &= f_{\text{ois}}(t) + \varphi_{\text{ois}}(t) + x(t), \\
r_{\text{proj}}(t) &= f_{\text{proj}}(t) + \varphi_{\text{proj}}(t) + x(t) + k(t).
\end{align}
The single shared \(x\) enforces that discounting and projection move together (a parallel shift in
rates moves both curves by the same amount), while \(k\) lets the basis fluctuate around its market
forward \(f_{\text{proj}}-f_{\text{ois}}\). The deterministic shifts \(\varphi\) make the model
reprice the initial curves exactly. Without them the model discount factor would be
\(P^{\text{mkt}}(T)\,\mathbb{E}[e^{-\int_0^T x}]\), and since \(\mathbb{E}[e^{-\int_0^T x}]>1\) the
model forward would sit below the market forward---in our data by \(10\bp\) at \(5\)Y and
\(22\)--\(37\bp\) at \(10\)Y. This is the Hull--White \(\varphi(t)\), or the CIR\(^{++}\)
shift~\cite{brigo2006}, generalised to stochastic variance, where it has no closed form; we
therefore fix it on the simulated path set itself (Section~\ref{sec:calib}). Because
\(\varphi\) is deterministic it never enters the factor dynamics. The factor dynamics under the pricing measure are
\begin{align}
\mathrm{d}x(t) &= -\lambda\, x(t)\dt + \sqrt{V(t)}\;\dW^{x}(t), \\
\mathrm{d}k(t) &= -\gamma\, k(t)\dt + \xi\;\dW^{k}(t), \\
V(t) &= w\,V_1(t) + (1-w)\,V_2(t), \\
\mathrm{d}V_i(t) &= \kappa_i\big(\theta_i(t) - V_i(t)\big)\dt + \varepsilon\,\sqrt{V_i(t)}\;\dW^{v_i}(t),
\quad i=1,2 .
\end{align}
The shared factor \(x\) carries a \emph{stochastic} instantaneous variance \(V\), built as a convex
mixture of two Cox--Ingersoll--Ross variances with a fast and a slow mean-reversion
(\(\kappa_1\equiv\kappa_f\), \(\kappa_2\equiv\kappa_s\)); this two-timescale construction is what
lets a single parameter set bend across maturities. Such a multi-factor CIR mixture is precisely a
low-order \emph{Markovian lift}: a finite-dimensional, affine, simulable process whose weighted sum
approximates a rough (or multi-scale) variance~\cite{gatheral2018,eleuch2019,abijaber2019multi,abijaber2019lifting,abijaber2019volterra,cuchiero2019}.
Truncating to two factors keeps the model tractable for large-scale Monte Carlo calibration while
retaining the two-timescale flexibility that a single CIR factor lacks. The Heston vol-of-vol
\(\varepsilon\) generates smile \emph{curvature}, and a term structure of leverage correlations
\(\rho(t)=\mathrm{corr}(\dW^{x},\dW^{v_1})\) generates the \emph{skew}. The long-run variance levels
\(\theta_1(t),\theta_2(t)\) and the correlation \(\rho(t)\) are represented as monotone
(PCHIP) splines through fixed nodes rather than constants. The Hull--White factor \(k\), with
constant volatility \(\xi\), carries the OIS-vs-projection basis, which in the EUR data is
\(+21\bp\) at \(1\)Y, \(+18\bp\) at \(5\)Y and \(+9\bp\) at \(10\)Y.

Setting \(\varepsilon=0\) collapses the variance to a deterministic term structure and the smile to
strike-flat; further flattening \(\theta\) and switching off \(k\) recovers the two- and
one-factor Gaussian Hull--White models. This nesting (Section~\ref{sec:ladder}) provides an
apples-to-apples decomposition of the fit.

A note on the name: ``Heston--Hull--White'' is also used for the equity--rates hybrid in which a
Heston stock is paired with Hull--White rates~\cite{grzelak2011}. Here there is no equity: the
Heston-type variance drives the \emph{rate} factor itself.

\subsection{A short rate on the projection curve}
Henrard~\cite[Sec.~5.3]{henrard2026} warns that short-rate models do not carry over to forward
curves: the projection curve's pseudo-discount factors are not prices of traded assets, are
arbitrary before the first fixing, and an asset would have to coincide
with the collateral curve. Our construction avoids the objection in two ways. First, the only
short rate that discounts is \(r_{\text{ois}}\), the collateral rate, which is the curve on which a
short-rate model is legitimate. Second, \(r_{\text{proj}}=r_{\text{ois}}+(f_{\text{proj}}-f_{\text{ois}})+k\)
is never used to discount; it enters only through the exponent of the projected rate over an
accrual period, i.e.\ as a continuously compounded spread over the collateral curve. Henrard shows
that this is the coherent way to represent a spread (his Sec.~5.2.5), because it
composes consistently across periods. \(\int r_{\text{proj}}\) is therefore a modelling device for
the Euribor fixing, not a claim that a Euribor bond exists. In the RFR markets the question does not
arise: the projection curve is the collateral curve, and \(k\) only adds noise around a zero basis.

\section{Instrument and volatility inversion}
\label{sec:instrument}
A caplet with accrual period \([t_{\text{fix}},t_{\text{fix}}+\tau]\) and strike \(K\) pays
\(\tau(R-K)^{+}\) at \(t_{\text{fix}}+\tau\), discounted on the collateral (OIS) curve. What
\(R\) is depends on the market, and the distinction matters for the volatility
(Section~\ref{sec:lit}).

\paragraph{EUR: forward-looking term rate.} 6M Euribor is fixed at the start of the period, so
\(R\) is known at \(t_{\text{fix}}\) and carries no volatility after it. In the model it is the
simple rate implied by the expected projection integral over the coming period,
\begin{equation}
R_{\text{term}} = \frac{1}{\tau}\Big(\exp\Big\{\mathbb{E}_{t_{\text{fix}}}\Big[\int_{t_{\text{fix}}}^{t_{\text{fix}}+\tau}
r_{\text{proj}}(u)\,\mathrm{d}u\Big]\Big\}-1\Big).
\end{equation}
Because \(x\) and \(k\) have linear drifts, the conditional expectation is exact whatever the
variance does after \(t_{\text{fix}}\):
\(\mathbb{E}_{t_{\text{fix}}}[\int r_{\text{proj}}] = \int f_{\text{proj}} + x(t_{\text{fix}})B_\lambda(\tau) +
k(t_{\text{fix}})B_\gamma(\tau)\), with \(B_a(\tau)=(1-e^{-a\tau})/a\) (in its exact
discrete-time form on the simulation grid). The Monte Carlo therefore needs only the state at the
fixing date.

\paragraph{USD and GBP: backward-looking compounded RFR.} SOFR and SONIA caplets pay the
overnight rate compounded through the accrual period~\cite{lyashenko2019}, which in
continuous time is
\begin{equation}
R_{\text{arr}} = \frac{1}{\tau}\Big(\exp\Big\{\int_{t_{\text{fix}}}^{t_{\text{fix}}+\tau}
r_{\text{proj}}(u)\,\mathrm{d}u\Big\}-1\Big),
\end{equation}
accrued over every business day of the period (Section~\ref{sec:calib}). In both markets the
forward of \(R\) is the simple forward off the projection curve, which is the forward the market
surface is stripped on; for the in-arrears rate this holds because, before its period starts, a
compounded overnight coupon has the same value as a term coupon on the same curve, whatever the
model.

\paragraph{Pricing and inversion.} Monte Carlo estimators of the discount factor
\(D=\mathbb{E}[e^{-\int_0^{t_{\text{fix}}+\tau} r_{\text{ois}}}]\), the forward
\(F=\mathbb{E}[R\,e^{-\int r_{\text{ois}}}]/D\) and the price
\(\mathrm{pv}=\mathbb{E}[(R-K)^{+}\tau\, e^{-\int r_{\text{ois}}}]\) are computed on the same path
set, and the price is inverted to a normal (Bachelier) volatility by
bisection on \(\sigma\):
\begin{equation}
\mathrm{pv} = D\,\tau\Big[(F-K)\,\Phi(d) + \sigma\sqrt{T_{\text{exp}}}\;\phi(d)\Big],
\qquad d=\frac{F-K}{\sigma\sqrt{T_{\text{exp}}}} .
\end{equation}
For the EUR term caplet \(T_{\text{exp}}=t_{\text{fix}}\). For the in-arrears caplets the rate
keeps setting until \(t_{\text{fix}}+\tau\); with volatility decaying linearly through the
period the effective expiry is
\(T_{\text{exp}}=t_{\text{fix}}+\tau/3\)~\cite{lyashenko2019,henrard2026}, and the first period
\([0,\tau]\), which is not yet fixed at trade date, is a live caplet. Model and market are
inverted with the same \(T_{\text{exp}}\), so they are compared on a common footing. This
matches how the data vendor prices these caps: Bloomberg's swap manager values a SOFR cap off the
same volatility cube as a strip of compounded in-arrears caplets, including the live first period,
in a one-factor Hull--White model, i.e.\ with volatility carried through the accrual period, of
which \(t_{\text{fix}}+\tau/3\) is the small-mean-reversion limit. We verified this treatment for SOFR; SONIA
caps follow the same compounded-in-arrears convention and are treated identically. The rate conventions (one-day
lookback, one rate cut-off day, two-day payment delay) shift dates by one
or two days. We ignore them; their convexity effect is negligible.

\section{Market data}
Three market inputs are fed in, all sourced from Bloomberg:
\begin{itemize}
  \item the OIS forward curve (\(360\) monthly points to \(30\)Y) for discounting;
  \item the 6M~Euribor projection forward curve (\(360\) monthly points) for the caplet underlying; and
  \item the caplet normal-volatility surface, stripped from the flat cap-vol quotes into
  per-period Bachelier volatilities indexed by \((t_{\text{fix}},K)\), with the expiry
  convention of Section~\ref{sec:instrument}.
\end{itemize}
The OIS curves are bootstrapped from OIS par rates and the Euribor curve is represented by
pseudo-discount factors. The Euribor curve is bootstrapped
from deposit, FRA and swap quotes in single-curve fashion; the fully consistent construction
discounts those instruments on €STR~\cite[Sec.~4.9]{henrard2026}, a refinement that
moves the Euribor forwards by a few basis points and that we leave to future work. Because the model
reprices whichever curves it is given (Section~\ref{sec:calib}), this affects the forwards, not the
volatility fit. For USD and GBP the projection curve is the OIS (SOFR, SONIA) curve itself: the
caplets reference the overnight rate, so the two curves coincide and only the EUR fit exercises the
basis factor. This
is also the vendor's own set-up: the USD RFR volatility cube uses the SOFR curve both to project
and to discount, and quotes cap volatilities in normal (Bachelier) terms.
The calibration set used here is the sub-surface with \(t_{\text{fix}}\le 10\)Y: \(20\) maturities
(\(0.5,1.0,\dots,10.0\)Y) \(\times\) \(17\) strikes (\(-1.5\%\) to \(+10\%\)), i.e.\ \(340\)
caplets, fitted \emph{jointly} with one parameter vector.

\section{Calibration}
\label{sec:calib}
\paragraph{Simulation.} Paths are generated on a GPU with an \(O(N)\)-memory Monte Carlo (running
integrals, never materialising the full path matrix for the pricing sums). The CIR variance
factors are sampled with the Andersen (2008) quadratic-exponential (QE)
scheme~\cite{andersen2008,broadie2006}. In the low-vol
regime the CIR is strongly Feller-violated (\(2\kappa\theta\ll\varepsilon^2\)); an Euler scheme
with reflection biases the variance upward---most severely at short maturities, where it inflated
the \(0.5\)Y ATM volatility from a market level of \(44\bp\) to over \(100\bp\)---whereas QE is
moment-exact and places an honest point mass at zero in the exponential regime.

\paragraph{Fitting the initial curves.} Only integrals of \(\varphi\) enter the prices, and only
over the caplet dates. For discounting,
\(\int_0^T\varphi_{\text{ois}} = \ln\mathbb{E}[e^{-\int_0^T x}]\), estimated on the same paths, so
that \(D(T)=P^{\text{mkt}}(T)\). For projection, \(\int_{t_{\text{fix}}}^{t_{\text{fix}}+\tau}
\varphi_{\text{proj}}\) is one constant per caplet in the exponent of \(R\), solved in closed form
so that \(F=\mathbb{E}[R\,e^{-\int r_{\text{ois}}}]/D\) equals the market simple forward. The
discount shift rescales price and \(D\) alike and leaves implied volatilities unchanged; the
projection shift places each strike at the right moneyness. On the daily grid the model forwards
then match the market to within \(0.2\bp\) at every maturity. In an HJM formulation the initial
curve is matched by construction; the shift \(\varphi\) is
what restores that property for a short-rate model.

\paragraph{Time step.} The rate is accrued on the simulation grid, so each accrual period must span
a whole number of steps. On the daily grid \(dt=1/252\) a \(6\)M period is exactly \(126\) steps and
a \(3\)M period \(63\), so every business day is represented, as in an RFR contract; at \(dt=0.02\) a
\(3\)M period would be \(12.5\) steps. In earlier runs the daily grid cut the USD ATM error from
\(5.2\) to \(3.6\bp\).

\paragraph{Objective.} Parameters \(\Theta\) are estimated by minimising a \emph{vega-weighted}
normal-volatility RMSE over the \(N\) caplets,
\begin{equation}
\mathcal{L}(\Theta)=\Big(\frac{1}{N}\sum_{j=1}^{N} w_j\,
\big(\sigma^{\text{mod}}_j(\Theta)-\sigma^{\text{mkt}}_j\big)^2\Big)^{1/2},
\qquad
w_j\propto\max\!\big(e^{-d_j^2/2},\,0.2\big)\,\big(1+\tfrac12\,\mathbf{1}_{\{t_{\text{fix},j}<1\}}\big),
\end{equation}
with \(d_j=(K^{\text{atm}}_j-K_j)\big/\big(\sigma^{\text{mkt}}_j\sqrt{\max(T_{\text{exp},j},0.1)}\big)\),
the weights normalised to mean one, and \(K^{\text{atm}}_j\) the strike of lowest market volatility
in caplet \(j\)'s maturity slice. The factor \(e^{-d^2/2}\) is the shape of the Bachelier vega:
it fits the liquid ATM tightly and down-weights the near-intrinsic deep-ITM and far-OTM wings, whose
quoted volatilities carry almost no vega and are numerically ill-defined; the \(20\%\) floor keeps
the wings in the fit, and the factor \(1.5\) on maturities below \(1\)Y counters the few short-dated
slices. All reported errors are unweighted. The parameters are optimised in a box-normalised space
\([0,1]^{27}\) by Adam~\cite{kingma2015} on the simultaneous-perturbation (SPSA)
gradient~\cite{spall1998}
\begin{equation}
\hat g(\Theta)=\frac{1}{M}\sum_{m=1}^{M}
\frac{\mathcal{L}(\Theta+c\,\Delta_m)-\mathcal{L}(\Theta-c\,\Delta_m)}{2c}\,\Delta_m ,
\end{equation}
with Rademacher directions \(\Delta_m\in\{\pm1\}^{27}\), \(M=8\), and cosine decay of the step
(\(3\times10^{-3}\)) and of the perturbation \(c\) (\(0.02\)) over \(800\) iterations. The estimator
is gradient-free, which suits the noisy, non-smooth Monte Carlo objective, and the \(2M+1\)
evaluations of an iteration are simulated in one batched pass. The loss converges by
\(\approx 150\) iterations and then oscillates at the SPSA noise floor; the running best is
retained.

\section{Results}
\subsection{Goodness of fit}
Table~\ref{tab:fit} reports the fit. The proposed model attains an RMSE of \(5.9\bp\), an ATM mean
absolute error of \(2.0\bp\) and a liquid-core (\(|K-F|<2\%\), \(F\) the market forward) RMSE of
\(2.8\bp\), as means \(\pm\) standard deviation over three Monte Carlo seeds at \(2\times10^5\) paths.
The remaining columns re-evaluate the two earlier calibrations and isolate each correction. The
term payoff (third to second column) lowers the RMSE by \(27\%\), entirely away from the money:
ATM and core errors are unchanged, while every maturity bucket from \(1\)Y improves. The curve fit
(second to first column) has the opposite effect: by measuring strikes against the market forward
it lowers the ATM and liquid-core errors by \(18\%\) and \(25\%\), with little change in the overall
RMSE. Figures~\ref{fig:surface3d} and~\ref{fig:review} show the fitted surface and its residuals.

\begin{figure}[!htbp]
\centering
\includegraphics[width=0.86\textwidth]{fig_eur_surface3d.png}
\caption{EUR caplet normal-volatility surface, \(t_{\text{fix}}\le10\)Y: model (coloured) and
market (black mesh). Model and market agree across the liquid part of the surface and separate only
in the short-dated deep in-the-money wing, where quotes are near intrinsic. Strikes are plotted on a
uniform index, with actual strikes on the axis labels.}
\label{fig:surface3d}
\end{figure}

\begin{table}[!htbp]
\centering
\begin{tabular}{lrrr}
\toprule
 & \textbf{proposed} & no curve fit & in-arrears \\
Metric (bp) & term payoff + curve fit & term payoff & average \\
\midrule
Overall normal-vol RMSE & \(\mathbf{5.90\pm0.02}\) & \(6.29\pm0.03\) & \(8.57\pm0.18\) \\
ATM mean absolute error & \(\mathbf{1.98\pm0.08}\) & \(2.42\pm0.01\) & \(2.34\pm0.01\) \\
Liquid core RMSE (\(|K-F|<2\%\)) & \(\mathbf{2.81\pm0.02}\) & \(3.73\pm0.01\) & \(3.66\pm0.01\) \\
\midrule
RMSE, \(0.5\)--\(1\)Y & \(13.7\) & \(14.7\) & \(16.8\) \\
RMSE, \(1\)--\(3\)Y & \(5.1\) & \(6.2\) & \(9.5\) \\
RMSE, \(3\)--\(7\)Y & \(4.5\) & \(4.2\) & \(7.2\) \\
RMSE, \(7\)--\(10\)Y & \(3.0\) & \(3.2\) & \(4.9\) \\
\bottomrule
\end{tabular}
\caption{Fit of the calibrated 2F lifted-CIR HWH model on the EUR surface (\(340\) caplets, \(27\)
parameters, daily step \(dt=1/252\), QE simulation, vega-weighted objective). Means \(\pm\) standard
deviation over three seeds at \(2\times10^5\) paths; all columns evaluated by the same procedure.}
\label{tab:fit}
\end{table}

\begin{figure}[!htbp]
\centering
\includegraphics[width=\textwidth]{fig_eur_review.png}
\caption{EUR calibration diagnostics. Top left: ATM term structure, market (black) and model (red).
Top right: smile slices at 1, 2, 5 and 10Y, market solid and model dashed. Bottom left: residual
\(\text{model}-\text{market}\) (bp) by maturity and strike. Bottom right: residual histogram for the
liquid core and the wings. Residuals are small from \(1\)Y; the main misfits are the \(0.5\)--\(1\)Y
ATM level (Section~\ref{sec:discussion}) and the near-intrinsic short-dated deep-ITM wing.}
\label{fig:review}
\end{figure}

\FloatBarrier
\subsection{Nested-model ladder}
\label{sec:ladder}
Table~\ref{tab:ladder} compares the nested restrictions, each calibrated and evaluated exactly as
the headline fit. The three models with \(\varepsilon=0\), whose smiles are strike-flat, reach
\(45\)--\(47\bp\) with ATM and liquid-core errors near \(20\bp\); the stochastic variance brings these
to \(5.9\), \(2.0\) and \(2.8\bp\). Neither a second Gaussian factor nor a deterministic volatility
term structure helps: stochastic volatility fits the surface at the money as well as in the wings.

\begin{table}[!htbp]
\centering
\begin{tabular}{lrrr}
\toprule
Model (bp) & RMSE & ATM MAE & core \\
\midrule
1F Gaussian HW         & 46.7 & 18.7 & 18.9 \\
2F Gaussian HW         & 47.1 & 18.4 & 18.2 \\
HW + deterministic vol & 45.3 & 21.2 & 21.3 \\
\textbf{HWH (proposed)} & \textbf{5.9} & \textbf{2.0} & \textbf{2.8} \\
\bottomrule
\end{tabular}
\caption{Nested-model ladder, EUR. All models fitted and evaluated as in Table~\ref{tab:fit}; the
first three have \(\varepsilon=0\) and strike-flat smiles. Seed standard deviation below
\(0.6\bp\) (RMSE) and \(0.1\bp\) (other columns).}
\label{tab:ladder}
\end{table}

\begin{figure}[!htbp]
\centering
\includegraphics[width=\textwidth]{fig_eur_ladder.png}
\caption{Nested-model ladder: RMSE bars (left), ATM term structure (centre), and the \(5\)Y smile
(right). The Gaussian/deterministic-vol rungs are strike-flat; only the stochastic-volatility model
follows the market smile, and only it tracks the ATM term structure. Same set-up as
Table~\ref{tab:ladder}.}
\label{fig:ladder}
\end{figure}

\FloatBarrier
\subsection{Calibrated parameters}
The model has \(27\) free parameters: two \(8\)-node long-run variance term structures
\(\theta_1,\theta_2\), a \(3\)-node leverage term structure \(\rho\), the two mean reversions
\(\kappa_f,\kappa_s\), the mixing weight \(w\), the vol-of-vol \(\varepsilon\), the level-factor and
basis-factor reversions \(\lambda,\gamma\), the basis vol \(\xi\), and the initial variance
\(V_0\). Tables~\ref{tab:theta} and \ref{tab:scalar} report the full calibrated set. The
\(\theta_i(t)\) and \(\rho(t)\) nodes are interpolated by monotone PCHIP splines; \(\theta\) is
reported as \(\sqrt{\theta}\times10^4\) (the instantaneous vol scale in bp, before mean reversion
pulls the caplet implied vol below it).

The estimates are economically sensible: a fast/slow variance pair (\(\kappa_f=1.56\),
\(\kappa_s=0.19\)) mixed with weight \(w=0.86\); a moderate vol-of-vol \(\varepsilon=0.10\); a
slowly mean-reverting level factor (\(\lambda=0.19\)); and a \emph{positive} leverage \(\rho\),
\(0.53\) at 1Y and \(0.63\) at 5Y before falling to \(0.18\) at 10Y, consistent with the right-skewed
(payer) cap smile. The basis factor is persistent (\(\gamma=0.03\)) with volatility
\(\xi=24\bp\) per \(\sqrt{\text{year}}\): the Euribor--OIS basis wanders slowly around its market
forward rather than snapping back to it. The fast reversion \(\kappa_f\) lies just above the lower
bound of its search range (\(1.5\)). Refitting with that bound relaxed to \(0.5\), warm-started from
the same parameters, moves \(\kappa_f\) to \(1.10\) with the other parameters essentially unchanged
and no gain in fit (RMSE \(5.92\) against \(5.90\bp\); ATM and liquid core slightly worse, \(2.6\) and
\(3.1\bp\)): \(\kappa_f\) is only weakly identified between about \(1\) and \(1.6\), and the bound does
not limit the fit.

\begin{table}[!htbp]
\centering
\begin{tabular}{lcccccccc}
\toprule
node (yr) & 0.25 & 0.5 & 1 & 2 & 3 & 5 & 7 & 10 \\
\midrule
\(\sqrt{\theta_1}\) (bp) & 82.7 & 143.7 & 151.5 & 160.5 & 141.1 & 176.3 & 167.5 & 180.1 \\
\(\sqrt{\theta_2}\) (bp) & 92.4 & 158.3 & 101.6 & 169.9 & 200.6 & 131.8 & 81.6 & 172.5 \\
\bottomrule
\end{tabular}
\caption{Calibrated long-run variance term structures (fast factor \(\theta_1\), slow factor
\(\theta_2\)), reported as \(\sqrt{\theta}\) in bp at the \(8\) spline nodes.}
\label{tab:theta}
\end{table}

\begin{table}[!htbp]
\centering
\begin{tabular}{lccccccccc}
\toprule
\(\kappa_f\) & \(\kappa_s\) & \(w\) & \(\varepsilon\) & \(\lambda\) & \(\gamma\) & \(\xi\) &
\(V_0\) & \\
\midrule
1.56 & 0.19 & 0.86 & 0.10 & 0.19 & 0.03 & \(2.4\!\times\!10^{-3}\) & 0 & \\
\midrule
\multicolumn{2}{l}{leverage} & \(\rho(1\text{Y})\) & \(\rho(5\text{Y})\) & \(\rho(10\text{Y})\) &
\multicolumn{4}{l}{} \\
& & 0.531 & 0.634 & 0.180 & \multicolumn{4}{l}{} \\
\bottomrule
\end{tabular}
\caption{Calibrated scalar parameters (top) and the \(3\)-node leverage term structure (bottom).}
\label{tab:scalar}
\end{table}

\FloatBarrier
\subsection{Cross-market robustness}
To check that the fit is not an EUR artefact, the identical pipeline (QE simulation, vega-weighted
objective, \(27\) parameters, \(800\) SPSA iterations, \(dt=1/252\)) was applied to the USD and GBP caplet
surfaces. USD and GBP are RFR markets (3M SOFR, 3M SONIA), priced as compounded in-arrears caplets;
EUR is the forward-looking IBOR (6M Euribor) case (Section~\ref{sec:instrument}). Table~\ref{tab:cross} reports the fit for
all three, with a single parameter set per market, despite the different underlyings, curve
conventions and surface sizes (USD/GBP carry \(697\) caplets on the 3M grid, including the first
period, vs \(340\) for EUR). The RFR headline is computed from \(t_{\text{fix}}\ge0.5\)Y: the first
caplet (\(t_{\text{fix}}=0\), effective expiry \(\tau/3\approx1\)M) has near-intrinsic wings whose
inverted volatilities are numerically meaningless, so we report it separately (RMSE \(63\bp\) USD,
\(77\bp\) GBP) rather than let it dominate the average. From \(3\)Y the three markets are fitted to similar accuracy (\(3\)--\(6\bp\) by
maturity bucket); below that the RFR markets are weaker, USD between \(1\) and \(3\)Y (\(12.9\bp\)) and
GBP below \(1\)Y (Section~\ref{sec:discussion}). On the GBP surface, which is the one whose curves are
unchanged, the same evaluation of the earlier calibration gives \(9.6\bp\) from \(6\)M, against
\(6.7\bp\) now.
Figure~\ref{fig:cross} shows the market-vs-model ATM term structure and \(5\)Y smile for each.

\begin{table}[!htbp]
\centering
\begin{tabular}{llrrrrr}
\toprule
market & payoff & caplets & RMSE & ATM MAE & liquid RMSE & full-surface RMSE \\
\midrule
EUR & 6M Euribor, term     & 340 & \(5.90\) & \(1.98\) & \(2.81\) & \(5.90\) \\
USD & 3M SOFR, in-arrears  & 697 & \(7.98\) & \(2.17\) & \(3.96\) & \(16.04\) \\
GBP & 3M SONIA, in-arrears & 697 & \(6.72\) & \(4.21\) & \(4.52\) & \(18.25\) \\
\bottomrule
\end{tabular}
\caption{Cross-market fit (bp), daily grid, curve fit included. Error columns over
\(t_{\text{fix}}\ge0.5\)Y; the last column includes the first RFR caplet. Seed standard deviation
of the RMSE \(\le0.27\bp\), below \(0.1\bp\) elsewhere.}
\label{tab:cross}
\end{table}

\begin{figure}[!htbp]
\centering
\includegraphics[width=\textwidth]{fig_cross_market.png}
\caption{Cross-market HWH calibration (QE, vega-weighted, daily grid \(dt=1/252\)): ATM term
structure (top row) and \(5\)Y smile (bottom row) for EUR, USD and GBP; market (black) vs.\ model
(red). The market term structures are piecewise flat because the stripped caplet volatility is
constant between quoted cap maturities.}
\label{fig:cross}
\end{figure}

\section{Discussion and further work}
\label{sec:discussion}
The model reproduces the EUR caplet surface across the whole strike/maturity range with a single
parameter set, and the RFR surfaces from \(6\)M out. Getting the instrument and the curves right
mattered as much as the model: treating Euribor as the forward-looking term rate it is removed a
quarter of the EUR error away from the money, and repricing the initial curves removed a quarter of
the liquid-core error. The remaining residuals are these.

\emph{The short end.} Below \(1.5\)Y the model undershoots the ATM level: in EUR \(39\) against
\(44\bp\) at \(0.5\)Y and \(61\) against \(72\bp\) at \(1\)Y, and in GBP \(35\) against \(52\bp\) at
\(0.5\)Y and \(52\) against \(78\bp\) at \(1\)Y; USD is closer (\(70\) against \(74\bp\) at
\(0.5\)Y). Part of this is in the data---the EUR \(1\)Y point sits above the \(1.5\)Y point, and the
stripped surface is piecewise flat between cap maturities---but a diffusion with mean-reverting
variance also struggles to build a short-dated hump. Jumps around policy dates~\cite{fontana2020disc,fontana2024overnight,andersenbang2024}
are the natural extension.

\emph{Near-intrinsic wings.} The short-dated deep-ITM and far-OTM strikes carry almost no vega; their
inverted volatilities are ill-conditioned, are down-weighted by design, and dominate the
full-surface RMSE of the RFR markets through the first caplet.

Planned extensions: (i) calibrating the full surface to \(29.5\)Y (the stripped surface already carries these
quotes); (ii) semi-analytic bond prices from Riccati equations, which would give \(\varphi\) without
simulation and which the next item needs anyway---the model is affine except for the leverage
term, whose covariance \(\rho\,\varepsilon\sqrt{V V_1}\) is not linear in the state unless \(w=1\),
so that term needs an approximation (e.g.\ linearising the degree-one homogeneous
\(\sqrt{V V_1}\) at the long-run variance ratio) checked against simulation; and (iii) an
out-of-sample check by repricing ATM
swaptions---which do not enter the calibration---with the calibrated parameters. Because a swaption probes the
\emph{correlation} between forward rates that a caplet surface only partly constrains, this is the
most informative remaining test; it requires bond reconstruction at the option expiry under
stochastic volatility (a conditional-Gaussian or nested-MC pricer), and is left as the natural next
step.

\clearpage
\appendix
\section{The quadratic-exponential (QE) variance scheme}
Each CIR variance factor is advanced over a step \(\Delta\) with the Andersen (2008) QE scheme.
Writing \(e=e^{-\kappa\Delta}\), the exact conditional mean and variance of \(V_{t+\Delta}\) given
\(V_t\) are
\begin{equation}
m = \theta + (V_t-\theta)\,e, \qquad
s^2 = \frac{V_t\,\varepsilon^2 e(1-e)}{\kappa} + \frac{\theta\,\varepsilon^2(1-e)^2}{2\kappa},
\qquad \psi = \frac{s^2}{m^2}.
\end{equation}
For \(\psi\le\psi_c\) (with \(\psi_c=1.5\)) a quadratic sample is drawn,
\begin{equation}
V_{t+\Delta} = a\,(b+Z)^2,\quad Z\sim N(0,1),\quad
b^2 = \tfrac{2}{\psi}-1+\sqrt{\tfrac{2}{\psi}\big(\tfrac{2}{\psi}-1\big)},\quad
a=\frac{m}{1+b^2};
\end{equation}
for \(\psi>\psi_c\) an exponential/point-mass sample is drawn,
\begin{equation}
V_{t+\Delta} = \begin{cases} 0, & U\le p,\\[2pt] \beta^{-1}\,\ln\dfrac{1-p}{1-U}, & U>p, \end{cases}
\qquad U\sim\mathrm{Unif}(0,1),\quad p=\frac{\psi-1}{\psi+1},\quad \beta=\frac{1-p}{m}.
\end{equation}
The scheme is exact in its first two moments and, in the exponential branch that dominates the
Feller-violated DM regime (\(2\kappa\theta\ll\varepsilon^2\)), places an honest atom at zero rather
than reflecting a negative Euler draw. The leverage correlation is imposed by driving the level
factor with the Brownian increment \(\rho\,\xi_1+\sqrt{1-\rho^2}\,Z_x\), where
\(\xi_1=(V_{1,t+\Delta}-m_1)/s_1\) is the standardised innovation of the fast variance factor (exact
in the quadratic branch, approximate in the exponential).

\clearpage
\begin{thebibliography}{99}
\small
\bibitem{abijaber2019multi} E.~Abi~Jaber and O.~El~Euch.
  Multifactor approximation of rough volatility models.
  \emph{SIAM Journal on Financial Mathematics}, 10(2):309--349, 2019. \doi{10.1137/18M1170236}
\bibitem{abijaber2019lifting} E.~Abi~Jaber.
  Lifting the Heston model.
  \emph{Quantitative Finance}, 19(12):1995--2013, 2019. \doi{10.1080/14697688.2019.1615113}
\bibitem{cuchiero2019} C.~Cuchiero and J.~Teichmann.
  Markovian lifts of positive semidefinite affine Volterra-type processes.
  \emph{Decisions in Economics and Finance}, 42(2):407--448, 2019. \doi{10.1007/s10203-019-00268-5}
\bibitem{gatheral2018} J.~Gatheral, T.~Jaisson, and M.~Rosenbaum.
  Volatility is rough.
  \emph{Quantitative Finance}, 18(6):933--949, 2018. \doi{10.1080/14697688.2017.1393551}
\bibitem{eleuch2019} O.~El~Euch and M.~Rosenbaum.
  The characteristic function of rough Heston models.
  \emph{Mathematical Finance}, 29(1):3--38, 2019. \doi{10.1111/mafi.12173}
\bibitem{abijaber2019volterra} E.~Abi~Jaber, M.~Larsson, and S.~Pulido.
  Affine Volterra processes.
  \emph{Annals of Applied Probability}, 29(5):3155--3200, 2019. \doi{10.1214/19-AAP1477}
\bibitem{vasicek1977} O.~Vasicek.
  An equilibrium characterization of the term structure.
  \emph{Journal of Financial Economics}, 5(2):177--188, 1977. \doi{10.1016/0304-405X(77)90016-2}
\bibitem{cir1985} J.~C.~Cox, J.~E.~Ingersoll, and S.~A.~Ross.
  A theory of the term structure of interest rates.
  \emph{Econometrica}, 53(2):385--407, 1985. \doi{10.2307/1911242}
\bibitem{hullwhite1990} J.~Hull and A.~White.
  Pricing interest-rate-derivative securities.
  \emph{Review of Financial Studies}, 3(4):573--592, 1990. \doi{10.1093/rfs/3.4.573}
\bibitem{heston1993} S.~L.~Heston.
  A closed-form solution for options with stochastic volatility with applications to bond and currency options.
  \emph{Review of Financial Studies}, 6(2):327--343, 1993. \doi{10.1093/rfs/6.2.327}
\bibitem{duffie2000} D.~Duffie, J.~Pan, and K.~Singleton.
  Transform analysis and asset pricing for affine jump-diffusions.
  \emph{Econometrica}, 68(6):1343--1376, 2000. \doi{10.1111/1468-0262.00164}
\bibitem{brigo2006} D.~Brigo and F.~Mercurio.
  \emph{Interest Rate Models---Theory and Practice: With Smile, Inflation and Credit}, 2nd edition.
  Springer Finance, Springer, 2006. \doi{10.1007/978-3-540-34604-3}
\bibitem{filipovic2009} D.~Filipovi\'{c}.
  \emph{Term-Structure Models: A Graduate Course}. Springer Finance, Springer, 2009.
  \doi{10.1007/978-3-540-68015-4}
\bibitem{duffiekan1996} D.~Duffie and R.~Kan.
  A yield-factor model of interest rates.
  \emph{Mathematical Finance}, 6(4):379--406, 1996. \doi{10.1111/j.1467-9965.1996.tb00123.x}
\bibitem{dfs2003} D.~Duffie, D.~Filipovi\'{c}, and W.~Schachermayer.
  Affine processes and applications in finance.
  \emph{Annals of Applied Probability}, 13(3):984--1053, 2003. \doi{10.1214/aoap/1060202833}
\bibitem{collindufresne2002} P.~Collin-Dufresne and R.~S.~Goldstein.
  Do bonds span the fixed income markets? Theory and evidence for unspanned stochastic volatility.
  \emph{Journal of Finance}, 57(4):1685--1730, 2002. \doi{10.1111/1540-6261.00475}
\bibitem{trolleschwartz2009} A.~B.~Trolle and E.~S.~Schwartz.
  A general stochastic volatility model for the pricing of interest rate derivatives.
  \emph{Review of Financial Studies}, 22(5):2007--2057, 2009. \doi{10.1093/rfs/hhn040}
\bibitem{andersenpiterbarg2010} L.~B.~G.~Andersen and V.~V.~Piterbarg.
  \emph{Interest Rate Modeling}, Volumes I--III. Atlantic Financial Press, 2010.
\bibitem{grzelak2011} L.~A.~Grzelak and C.~W.~Oosterlee.
  On the Heston model with stochastic interest rates.
  \emph{SIAM Journal on Financial Mathematics}, 2(1):255--286, 2011. \doi{10.1137/090756119}
\bibitem{bgm1997} A.~Brace, D.~Gatarek, and M.~Musiela.
  The market model of interest rate dynamics.
  \emph{Mathematical Finance}, 7(2):127--155, 1997. \doi{10.1111/1467-9965.00028}
\bibitem{hagan2002} P.~S.~Hagan, D.~Kumar, A.~S.~Lesniewski, and D.~E.~Woodward.
  Managing smile risk.
  \emph{Wilmott Magazine}, September 2002, 84--108.
\bibitem{bianchetti2010} M.~Bianchetti.
  Two curves, one price.
  \emph{Risk}, 23(8):74--80, 2010. Full version: \doi{10.2139/ssrn.1334356}
\bibitem{mercurio2010} F.~Mercurio.
  Interest rates and the credit crunch: new formulas and market models.
  Bloomberg Portfolio Research Paper No.\ 2010-01-FRONTIERS, 2010.
  \url{https://ssrn.com/abstract=1332205}
\bibitem{mercurioxie2012} F.~Mercurio and Z.~Xie.
  The basis goes stochastic.
  \emph{Risk}, 25(12):78--83, 2012.
\bibitem{grbac2015} Z.~Grbac and W.~J.~Runggaldier.
  \emph{Interest Rate Modeling: Post-Crisis Challenges and Approaches}.
  SpringerBriefs in Quantitative Finance, Springer, 2015. \doi{10.1007/978-3-319-25385-5}
\bibitem{fontana2016hjm} C.~Cuchiero, C.~Fontana, and A.~Gnoatto.
  A general {HJM} framework for multiple yield curve modelling.
  \emph{Finance and Stochastics}, 20(2):267--320, 2016. \doi{10.1007/s00780-016-0291-5}
\bibitem{cfg2019} C.~Cuchiero, C.~Fontana, and A.~Gnoatto.
  Affine multiple yield curve models.
  \emph{Mathematical Finance}, 29(2):568--611, 2019. \doi{10.1111/mafi.12183}
\bibitem{fontana2020disc} C.~Fontana, Z.~Grbac, S.~G\"umbel, and T.~Schmidt.
  Term structure modelling for multiple curves with stochastic discontinuities.
  \emph{Finance and Stochastics}, 24(2):465--511, 2020. \doi{10.1007/s00780-020-00416-5}
\bibitem{fontana2021cbi} C.~Fontana, A.~Gnoatto, and G.~Szulda.
  Multiple yield curve modelling with {CBI} processes.
  \emph{Mathematics and Financial Economics}, 15(3):579--610, 2021. \doi{10.1007/s11579-020-00289-4}
\bibitem{fontana2023rollover} C.~Fontana, S.~Pavarana, and W.~J.~Runggaldier.
  A stochastic control perspective on term structure models with roll-over risk.
  \emph{Finance and Stochastics}, 27(4):903--932, 2023. \doi{10.1007/s00780-023-00515-z}
\bibitem{fontana2024overnight} C.~Fontana, Z.~Grbac, and T.~Schmidt.
  Term structure modeling with overnight rates beyond stochastic continuity.
  \emph{Mathematical Finance}, 34(1):151--189, 2024. \doi{10.1111/mafi.12415}
\bibitem{fontana2025geometry} C.~Fontana, G.~Lanaro, and A.~Murgoci.
  The geometry of multi-curve interest rate models.
  \emph{Quantitative Finance}, 25(2):323--342, 2025. \doi{10.1080/14697688.2024.2409276}
\bibitem{henrard2026} M.~Henrard.
  \emph{Interest Rate Modelling in the Multi-Curve Framework: Foundations, Evolution, Transition,
  and Implementation}, 2nd edition. Applied Quantitative Finance, Palgrave Macmillan, 2026.
  \doi{10.1007/978-3-032-02685-9}
\bibitem{lyashenko2019} A.~Lyashenko and F.~Mercurio.
  Looking forward to backward-looking rates: a modeling framework for term rates replacing LIBOR.
  SSRN Working Paper 3330240, 2019. \doi{10.2139/ssrn.3330240}
\bibitem{fontana2023caplet} C.~Fontana.
  Caplet pricing in affine models for alternative risk-free rates.
  \emph{SIAM Journal on Financial Mathematics}, 14(1):SC1--SC16, 2023. \doi{10.1137/22M1513691}
\bibitem{andersenbang2024} L.~Andersen and D.~Bang.
  Spike and hike modeling for interest rate derivatives: with an application to SOFR caplets.
  \emph{Quantitative Finance}, 24(8):1017--1033, 2024. \doi{10.1080/14697688.2024.2364800}
\bibitem{broadie2006} M.~Broadie and \"{O}.~Kaya.
  Exact simulation of stochastic volatility and other affine jump diffusion processes.
  \emph{Operations Research}, 54(2):217--231, 2006. \doi{10.1287/opre.1050.0247}
\bibitem{andersen2008} L.~B.~G.~Andersen.
  Simple and efficient simulation of the Heston stochastic volatility model.
  \emph{Journal of Computational Finance}, 11(3):1--42, 2008. \doi{10.21314/JCF.2008.189}
\bibitem{spall1998} J.~C.~Spall.
  Implementation of the simultaneous perturbation algorithm for stochastic optimization.
  \emph{IEEE Transactions on Aerospace and Electronic Systems}, 34(3):817--823, 1998.
  \doi{10.1109/7.705889}
\bibitem{kingma2015} D.~P.~Kingma and J.~Ba.
  Adam: a method for stochastic optimization.
  \emph{International Conference on Learning Representations (ICLR)}, 2015.
  \doi{10.48550/arXiv.1412.6980}
\end{thebibliography}

\end{document}
